\section*{General binary $k=3$ class: corollary of M14--M15}
\label{sec:m16-general-binary-class}

Fix $0<\eta\le1$ and let $\mathcal A_{3,\mathrm{bin}}$ denote the
finite-dimensional real or complex projected-root class with unit leaves,
orthogonal projectors, three binary internal nodes, and no law-sharing
constraint: each internal vertex may use an arbitrary bilinear law satisfying
the common operator-norm and closure-defect budgets $M$ and
$\rho=\eta M$.  The two binary topologies are the ordered chain and the
ordered branching tree.

\textbf{Corollary.}  For either topology, the class-level normalized constant
is
\[
  \boxed{
  C_{3,\mathcal A_{\mathrm{bin}},\mathrm{chain}}^P(\eta)
  =C_{3,\mathcal A_{\mathrm{bin}},\mathrm{branch}}^P(\eta)
  =W_3(\eta),
  }
\]
where
\[
 W_3(\eta)=
 \begin{cases}
 \sqrt{4-3\eta^2}, & 0<\eta\le\sqrt{2/3},\\[1mm]
 2/(\sqrt3\,\eta), & \sqrt{2/3}\le\eta\le1.
 \end{cases}
\]

\textbf{Proof.}  The phrase ``no law-sharing constraint'' is exactly the
admissible-law convention used by M14 and M15: the laws at distinct vertices
may be selected independently, with no restriction on their individual
coefficients beyond the common budgets.  Thus the chain component of
$\mathcal A_{3,\mathrm{bin}}$ is the M14 class, and the branching component
is the M15 class.  M14 supplies the chain upper bound and a dimension-two
attaining witness; M15 supplies the corresponding branching upper bound and
witness.  Equality of the two class-level suprema with $W_3$ follows.
\hfill$\square$

\textbf{Boundary.}  This corollary does not determine the smaller supremum
under a repeated same-law constraint, the broader gated-planar repeated-law
subclass, higher-arity or arbitrary finite topologies, or growing-tree
families.
